% kernel/engineering_objects.tex

\chapter{Engineering Objects}
\label{chap:engineering-objects}

Engineering activity concerns entities that possess enduring meaning
within the engineering domain. The Knowledge Kernel calls these
entities engineering objects.

\section{Engineering Object}

\begin{definition}[Engineering object]
\label{def:engineering-object}
An engineering object is a durable engineering entity whose identity is
independent of its representations, observations, realizations, and
computational models.
\end{definition}

Durability does not imply immutability. An engineering object may be
created, modified, related to other objects, observed under changing
conditions, and represented in many forms while retaining one
engineering identity.

An engineering object is therefore not defined by a particular file,
coordinate system, database record, visual geometry, or software class.
Those may carry information about the object, but none of them alone is
the object.

\section{Object and Description}

The distinction between an engineering object and its descriptions is
fundamental. A state describes the object under a specified condition. A
representation communicates information about it. An observation
records evidence concerning it. A realization interprets or embodies it
within a metric or physical context.

These descriptions may change independently. A new representation does
not create a new engineering object. A revised observation does not by
itself change the object. A different metric realization does not alter
its identity.

\section{Persistence}

Engineering objects persist across workflows, projects, tools, and
representations. Persistence requires an identity that can be referred
to independently of transient processing states.

This distinction supports exchange, comparison, revision, provenance,
and long-term engineering responsibility. It also prevents temporary
preview geometry, visualization primitives, or runtime selections from
being mistaken for durable engineering objects.

\section{Relations}

Engineering objects may participate in explicit relations. Relations
express engineering meaning between objects without collapsing their
identities.

An alignment may use a profile, intersect a platform edge, participate
in a project, or be observed by a measurement campaign. These relations
connect engineering objects while preserving the responsibility and
identity of each participant.

\section{Canonical Interpretation}

\begin{proposition}
An engineering object remains distinguishable from every state,
representation, observation, realization, and computational model that
refers to it.
\end{proposition}

\begin{proposition}
Changes of representation or realization do not by themselves create a
new engineering object.
\end{proposition}

\begin{proposition}
Engineering relations connect objects without merging their identities.
\end{proposition}

The general engineering-object concept provides the basis for the later
identification of alignments and other railway entities as durable
objects within the engineering domain.
